Для целей принятия решения относительно окончательного выбора пути профессиональной самореализации составлена таблица с указанием наиболее значимых характеристик будущего трудоустройства (\cref{tab:comparison}). 

При сравнении учитывались расчетные значения чистого потенциального дохода в месяц: для первого пути (найм) показатель составил $104400$ руб., для второго пути (самозанятость) — от $80895$ до $82743$ руб. при сопоставимой временной нагрузке.

\begin{longtblr}[
    caption = {Сравнение вариантов будущего трудоустройства},
    label = {tab:comparison},
]{
    colspec = {X[1,l,m] X[1,l,m] X[1,l,m] X[1,l,m]},
    hlines, vlines,
    row{1,2} = {bg=gray9, c},
    rowhead = 1,
}
    \SetCell[c=2]{c} \textit{Путь 1 (Трудоустройство в компанию)} & & \SetCell[c=2]{c} \textit{Путь 2 (Самозанятость / Фриланс)} & \\
    \textit{Преимущества} & \textit{Недостатки} & \textit{Преимущества} & \textit{Недостатки} \\
    Более высокий и стабильный уровень дохода ($104400$ руб./мес). Наличие социального пакета. & Жесткая привязка к рабочему графику и локации (офис/производство). & Гибкость графика и возможность удаленной работы. & Финансовая нестабильность, отсутствие социальных гарантий. \\
    Возможность профессионального нетворкинга, работа в команде. & Подчинение корпоративным регламентам, ограниченный выбор проектов. & Самостоятельный выбор заказчиков и стека технологий. & Изоляция, дефицит живого общения. Необходимость самостоятельно искать клиентов. \\
\end{longtblr}

На основании проведенного анализа в качестве приоритетной стратегии профессионального развития выбран Путь 1 --- \textit{трудоустройство в компанию}. Данный выбор обусловлен личностными характеристиками (экстраверсия, потребность в коллективной работе и коммуникации), а также стремлением к финансовой стабильности и участию в крупных, ресурсоемких проектах, недоступных для индивидуальных исполнителей. Фриланс отвергнут по причине нестабильности потока заказов и негативного влияния удаленной работы на личную продуктивность.

Для компенсации выявленных недостатков выбранной стратегии предлагаются следующие меры:
\begin{itemize}
    \item \textit{Нивелирование жесткости графика:} применение методов тайм-менеджмента для оптимизации рабочих процессов и обсуждение с работодателем возможности гибридного формата работы (например, 1--2 дня в неделю удаленно).
    \item \textit{Компенсация рутины и ограниченности проектов:} инициация собственных исследовательских или оптимизационных задач внутри компании, а также ведение личных open-source проектов в свободное время для удовлетворения потребности в творческой реализации.
\end{itemize}